\exnumber{01}
\extitle{UX VII - Ex 2 - Empathie}
\exfiles{ex01/archive zip, nomme Name\_FirstName\_DeliverableName\_Date\_VersionNumber}
\exauthorize{Ceux nommés dans l'exercice}
\makeheaderfiles

\textbf{Instructions :}

Utilisez au moins un des outils Empathy du processus de design.

\textbf{Livrables :}

\begin{itemize}
\item Un document PDF avec le script du contexte de l'outil
\item Outil complété (scripts des entretiens utilisateurs, résultats d'une analyse quantitative, etc.)
\item Un document PDF avec les apprentissages clés de cette phase
\item Déposez les documents sur Github, dans un zip nommé « ProjectTitle\_Name\_FirstName », et chaque livrable ainsi : « Name\_FirstName\_DeliverableName\_Date\_VersionNumber ». Exemple : « Doe\_John\_MidFiPrototype\_05242024\_V1 ».
\item Ajoutez un titre dans tous les documents écrits.
\item Soignez l'aspect global des livrables. Ils doivent être clairs, simples et lisibles. Vous pouvez regarder canva.com, sans surcharger le dessin de détails.
\end{itemize}

\textbf{Pistes :}

\begin{itemize}
\item Utilisez autant d'outils que nécessaire, et assurez-vous de faire tout le processus en 4 jours ouvrés !
\end{itemize}

\textbf{Critères :}

\begin{itemize}
\item Entretiens utilisateurs
\begin{itemize}
\item Le script est le guide d'entretien. Assurez-vous qu'il y a :
\begin{itemize}
\item l'objectif principal de l'entretien
\item 1 à 3 hypothèses
\item une introduction appropriée qui décrit l'intervieweur et l'objectif de cet entretien
\item 3 à 5 sujets principaux
\item quelques questions de screening (données démographiques)
\item questions ouvertes, max 10
\item un processus en entonnoir (>)
\item une conclusion avec les prochaines étapes et un remerciement.
\end{itemize}
\item L'outil complété est soit :
\begin{itemize}
\item 5 à 10 scripts d'utilisateurs
\item ou un diagramme d'affinité \url{https://www.youtube.com/watch?v=VngWHIE4k9s} avec les éléments principaux
\end{itemize}
\end{itemize}
\end{itemize}

\begin{itemize}
\item Analytique
\begin{itemize}
\item Le script est le formulaire. Assurez-vous qu'il y a :
\begin{itemize}
\item L'objectif principal du formulaire
\item 1 à 3 hypothèses
\item Une introduction correcte qui décrit l'intervieweur et le but de cet entretien
\item 3 à 5 sujets principaux
\item Quelques questions de filtrage (données démographiques)
\item Questions fermées, sans biais, 10 max
\item Un processus en entonnoir (>)
\item Une conclusion avec les prochaines étapes et un remerciement.
\item Une stratégie de diffusion
\end{itemize}
\item L'outil complété est une série de données quantitatives à partir d'au moins 100 réponses.
\end{itemize}
\end{itemize}

\begin{itemize}
\item Analyse concurrentielle
\begin{itemize}
\item Le script est la stratégie. L'étudiant explique ce qu'il cherche et comment il prévoit de le trouver.
\item L'outil complété est un rapport. Assurez-vous qu'il y a ;
\begin{itemize}
\item Il y a une introduction
\item Il y a un paysage concurrentiel \url{https://venngage-wordpress.s3.amazonaws.com/uploads/2019/03/g2crowd-grid.png}
\item Avec 2 axes pertinents
\item 10 organisations apparaissent sur le paysage concurrentiel
\item Il y a un tableau concurrentiel \url{https://venngage-wordpress.s3.amazonaws.com/uploads/2019/03/product-features-comparison-matrix.png}
\item Avec 5 à 10 critères
\item 10 organisations apparaissent dans le tableau
\end{itemize}
\end{itemize}
\end{itemize}

